\chapter{Limiti noti}
\label{ch:limiti}

Questo capitolo raccoglie i limiti strutturali dei modelli, documentati esplicitamente perché
incidono sull'interpretazione dei risultati.

\section{Apparecchi sempre accesi sovradimensionati}

Come anticipato nel paragrafo~\ref{sec:m3}, il vincolo di attività permanente
\eqref{eq:always-on} è un'\textbf{uguaglianza}: gli apparecchi sempre accesi devono assorbire in
ogni istante almeno $\sum_j p^{\text{ao}}_j (1-\delta)$ watt.

Se il questionario dichiara apparecchi la cui somma supera il consumo medio effettivo
dell'abitazione, il modello sovrastima per costruzione, indipendentemente dalla qualità della
soluzione trovata.

\begin{center}
\begin{tabular}{@{}l r r l@{}}
\toprule
\textbf{Abitazione} & \textbf{Segnale medio} & \textbf{Always-on dichiarati} & \textbf{Esito} \\
\midrule
A & $238.6\unit{W}$ & $125.0\unit{W}$ & regolare \\
B & $186.5\unit{W}$ & $125.0\unit{W}$ & regolare \\
C & $\mathbf{159.0}\unit{\mathbf{W}}$ & $\mathbf{325.0}\unit{\mathbf{W}}$ & \textbf{sovrastima strutturale} \\
D & $757.8\unit{W}$ & $125.0\unit{W}$ & regolare \\
E & $199.1\unit{W}$ & $125.0\unit{W}$ & regolare \\
F & $715.1\unit{W}$ & $225.0\unit{W}$ & regolare \\
\bottomrule
\end{tabular}
\end{center}

Nell'abitazione C il questionario dichiara due frigoriferi e un congelatore, per un totale di
$325\unit{W}$ tipici, a fronte di un consumo medio misurato di $159\unit{W}$. L'errore
sull'energia totale risulta del $119\%$ e il residuo medio è negativo (circa
$-165\unit{W}$): il modello attribuisce più energia di quanta ne sia stata misurata.

\paragraph{Motivazione.} Non è un problema del solver né dei coefficienti. O il questionario dichiara
apparecchi che di fatto non consumano, o i valori di potenza tipica della base di conoscenza sono
troppo alti per quegli specifici modelli.

\paragraph{Possibile correzione.} Rilassare il vincolo~\eqref{eq:always-on} da uguaglianza a
disuguaglianza ($\le 1$), permettendo lo stato spento quando il segnale non sostiene il consumo
dichiarato, eventualmente con una penalità per lo spegnimento. Finché il vincolo resta
un'uguaglianza quell'abitazione non è utilizzabile.

\section{Allocazione greedy della quota settimanale}

Poiché i giorni sono risolti separatamente, la quota settimanale viene consumata in ordine
cronologico senza previsione: i primi giorni della settimana la spendono liberamente e i
successivi ne subiscono le conseguenze.

Una distribuzione ottima richiederebbe di risolvere i sette giorni congiuntamente, moltiplicando
per sette la dimensione del problema. Il compromesso adottato evita comunque l'esito implausibile
in cui un apparecchio dichiarato bisettimanale risulta acceso tutti i giorni.

\section{Quota settimanale su blocchi di cinque giorni}

I blocchi di analisi durano cinque giorni e vengono elaborati indipendentemente, quindi il
contatore settimanale riparte a ogni blocco. Un blocco di cinque giorni riceve così l'intera quota
settimanale del questionario, il che è \textbf{leggermente permissivo}: un apparecchio dichiarato
due volte a settimana ne riceve due in cinque giorni anziché circa $1.4$ in proporzione.

Inoltre l'azzeramento interno al blocco, previsto al cambio di settimana ISO, si verifica solo nei
blocchi che contengono un lunedì.

\section{Assenza di verità di riferimento}

Il limite più importante non riguarda la formulazione ma la validazione. Non esistendo etichette
temporali, \textbf{non è possibile misurare l'accuratezza dell'attribuzione}. Le metriche
disponibili sono indirette:

\begin{center}
\begin{tabularx}{\textwidth}{@{}l X@{}}
\toprule
\textbf{Metrica} & \textbf{Che cosa misura, e che cosa non misura} \\
\midrule
Errore di ricostruzione
& Quanto la somma degli apparecchi stimati aderisce al segnale. \textbf{Non} dice se
  l'attribuzione fra apparecchi è corretta: una soluzione che assegna tutto all'apparecchio
  sbagliato può avere errore nullo. \\
\addlinespace
Errore sull'energia totale
& Scostamento fra energia attribuita e misurata sul periodo. Stesso limite. \\
\addlinespace
Coerenza temporale
& Percentuale di cicli con durata plausibile. Indiretta ma indipendente dal segnale. \\
\addlinespace
Plausibilità rispetto al questionario
& Il criterio più informativo di fatto disponibile: verificare che gli apparecchi ricevano
  consumo negli orari e nei mesi dichiarati. È esattamente ciò che i prior descritti in questo
  documento impongono, quindi non costituisce una validazione indipendente. \\
\bottomrule
\end{tabularx}
\end{center}

\begin{warnbox}{Implicazione per il committente}
I risultati vanno letti come una \textbf{ripartizione plausibile} dei consumi, coerente con quanto
dichiarato dall'utente e con il segnale misurato, non come una misura verificata del consumo per
apparecchio. Una validazione vera richiederebbe un sottoinsieme di abitazioni strumentate con
misuratori a livello di presa, da usare come riferimento.
\end{warnbox}

\section{Tempo di calcolo e ottimalità}

Il budget di 60 secondi per giorno non garantisce di raggiungere l'ottimo. Sui modelli $L_1$
risolti con HiGHS il limite viene generalmente rispettato con ampio margine, ma la difficoltà
cresce rapidamente con il numero di apparecchi e con la finezza della discretizzazione.

Quando il budget si esaurisce il solver restituisce la migliore soluzione trovata fino a quel
momento, che può essere lontana dall'ottimo. In quel caso il confronto fra modelli diversi perde
significato, perché si finirebbe per confrontare la qualità della ricerca anziché la formulazione.
Le varianti a 30 minuti dimezzano il numero di variabili e sono la prima leva da usare quando il
tempo di calcolo diventa critico.
